\subsection{A characterization of Hilbert space}\label{subsec:hilbert-preliminaries-recognition}
\paraid{p-a8af9823b28c}
We use the following characterization of Hilbert space.
It is Toru\'nczyk's theorem
\cite[p.~248, condition~(i)]{Torunczyk1981},
with its correction
\cite[Section~C]{Torunczyk1985}.

\begin{theorem}[{Toru\'nczyk, \cite[p.~248, condition~(i)]{Torunczyk1981}, \cite[Section~C]{Torunczyk1985}}]\label{thm:torunczyk}
\paraid{p-e19e6ccab7f7}
Let
$\MetricSpace{\RecognitionTarget }{\TargetMetric }$
be a nonempty complete separable metric AR,
and let
$\HilbertCube =\UnitInterval^{\NonnegativeIntegers}$
be the Hilbert cube.
Equip
$\NonnegativeIntegers$
with the discrete topology.
Then
$\RecognitionTarget $
is homeomorphic to the real Hilbert space
$\SeparableHilbertSpace$
if and only if,
for every continuous map
$\InputMap \colon \NonnegativeIntegers\times \HilbertCube \to \RecognitionTarget $
and every open cover
$\OpenCover$
of
$\RecognitionTarget $,
there exists a continuous map
$\ApproximatingMap \colon \NonnegativeIntegers\times \HilbertCube \to \RecognitionTarget $
that is
$\OpenCover$-close to
$\InputMap$
and whose image family
$\{\ApproximatingMap (\{\FamilyIndex \}\times \HilbertCube )\}_{\FamilyIndex \in\NonnegativeIntegers}$
is discrete in
$\RecognitionTarget $.
\end{theorem}
